\documentclass[10pt,a4paper]{amsart}
\usepackage[margin=19mm]{geometry}
\usepackage{amsmath,amssymb,mathtools}
\usepackage[scr=rsfs,cal=euler]{mathalfa}
\usepackage{xfrac}
\usepackage[unicode,hidelinks,bookmarks=false]{hyperref}
\newcommand{\ie}{i.e.\ }
\newcommand{\cf}{cf.\ }
\newcommand{\resp}{respectively\ }
\newcommand{\Subsection}{\subsection}
\providecommand{\nobreakdash}{\nobreak\hbox{-}}
\setlength{\emergencystretch}{2em}
\multlinegap0pt
\theoremstyle{plain}
\newtheorem{theorem}{Theorem}
\newtheorem{lemma}[theorem]{Lemma}
\newtheorem{propo}[theorem]{Proposition}
\newtheorem{corollary}[theorem]{Corollary}
\theoremstyle{definition}
\newtheorem{definition}[theorem]{Definition}
\theoremstyle{remark}
\newtheorem{remark}[theorem]{Remark}
\title[A Unified Kantorovich Duality for Multimarginal Optimal Transport]{A Unified Kantorovich Duality for Multimarginal Optimal Transport: the compact-case Kantorovich duality (Theorem 1)}
\author{Yehya Cheryala, Mokhtar Z. Alaya, Salim Bouzebda}
\date{Source: arXiv:2601.17171v2 (CC BY 4.0)}
\begin{document}
\maketitle

\noindent\emph{Source and licence.} A Unified Kantorovich Duality for Multimarginal Optimal Transport. Version arXiv:2601.17171v2 (CC BY 4.0), distributed under the Creative Commons Attribution 4.0 International licence, http://creativecommons.org/licenses/by/4.0/, as recorded in the arXiv licence metadata for this exact version and shown on the abstract page https://arxiv.org/abs/2601.17171v2. Full licence notice: licenses/arxiv-2601.17171-CC-BY-4.0.txt.\par\smallskip

\noindent\textit{Editorial scope.} Counted unit: the article's Theorem 1 (source0.tex lines 380--389, source label \texttt{KANTO}), the Kantorovich duality formula for multimarginal optimal transport with compact metric spaces and continuous nonnegative cost, $\mathsf{MOT}_c(\mu_1,\ldots,\mu_K)=\inf_\pi I(\pi)=\sup_{\boldsymbol f\in\mathcal{F}_c}J(\boldsymbol f)$, delivered with the article's complete original proof (source0.tex lines 1743--1889; the article prints this proof in its appendix as the run-in block headed ``Proof of Theorem~\ref{KANTO}'', ending with the article's own $\Box$, and the block is delivered here exactly from its heading line through that closing symbol, with the blank lines that separate it from the next appendix heading excluded). No mathematics is written, repaired or re-derived: statement and proof are the article's own text line for line, in the article's own notation ($\mathcal{P}(X_k)$, $\Pi(\mu_1,\dots,\mu_K)$, $I(\pi)$, $\mathcal{F}_c$, $J(\boldsymbol f)$, $\Theta$, $E=\mathcal{C}(X_1\times\cdots\times X_K)$, $\mathcal{M}(X_1\times\cdots\times X_K)$), and no line of the counted statement or of the counted proof is omitted, substituted or re-flowed. Required context retained uncounted: the article's Lemma 3.1, the weak duality relation (source0.tex lines 331--345, source label \texttt{easyy}), which the last step of the counted proof invokes; and the article's Theorem A.6, its restatement of the Riesz representation theorem (source0.tex lines 3040--3055, source label \texttt{RIEZ}), which the counted proof invokes in its first sentence; that theorem's heading carries the article's citation \texttt{knapp2007}, whose bibliography entry is transcribed in the wrapper's bibliography. Editorial formatting only, in addition: an AMS \texttt{amsart} preamble that declares the theorem environments the retained lines use; the article's own class and style files (\texttt{jmlr.sty} and \texttt{macros.sty}) are neither shipped nor input, and the retained lines use only standard AMS commands, so no macro from those files had to be restated; sequential numbering of the retained environments in order of appearance, so that the retained lemma, theorem and appendix theorem are numbered 1, 2 and 3, whereas the article prints them as Lemma 3.1, Theorem 1 and Theorem A.6; the article's equation numbering is not reproduced, so displayed formulas inside the retained spans are unnumbered or numbered consecutively, and every cross-reference inside the retained text resolves to the wrapper's numbering. Exactly one bibliography entry is delivered, the cited key \texttt{knapp2007}, transcribed from the article's own \texttt{thebibliography} entry at source0.tex lines 3221--3225 with that entry's \texttt{\textbackslash newblock} line breaks removed and the article's printed label 19 kept; no other key is cited inside a retained span. No asset-inclusion command, no drawing environment and no image asset occurs in this package.

\section*{Retained context: weak duality and the Riesz representation theorem}

\begin{lemma}\label{easyy}
Under the above assumptions on the topological spaces \((X_k)_{k=1}^K\)
endowed with their Borel \(\sigma\)-algebras, the marginals
\((\mu_k)_{k=1}^K\), and the bounded measurable cost function
\(c:X_1\times\cdots\times X_K\to\mathbb R\), the following weak duality relation holds:
\begin{align*}
\sup_{\boldsymbol f\in\mathcal F_c \cap \prod_{k=1}^K \mathcal C_b(X_k)}
J(\boldsymbol f)
\;\le\;
\sup_{\boldsymbol f\in\mathcal F_c}
J(\boldsymbol f)
\;\le\;
\inf_{\pi\in\Pi(\mu_1,\dots,\mu_K)} I(\pi).
\end{align*}
\end{lemma}


\section*{Retained context: the appendix statement of the Riesz representation theorem}

\begin{theorem}[Riesz representation theorem \cite{knapp2007}]\label{RIEZ}
Let $X_1, \dots, X_K$ be compact metric spaces and denote by $ E=\mathcal{C}(X_1 \times\ldots \times X_K)$ the space of continuous and real-valued functions
on the product, equipped with the supremum norm. 
Then, the dual space $E^*$ can be identified with the space $\mathcal{M}\!(X_1 \times\ldots \times X_K)$ of finite signed regular Borel measures, equipped with the total variation norm. The correspondence is given by
\begin{align*}
\mu \in \mathcal M(X_1 \times \cdots \times X_K)
\;\longmapsto\;
\mathcal{L}_\mu \in E^{*},
\end{align*}
where, for every $u \in E$,
\begin{align*}
\mathcal{L}_\mu(u) 
= \int_{X_1 \times\ldots \times X_K} 
u(x_1, \dots, x_K)\, d\mu(x_1, \dots, x_K).
\end{align*}
\end{theorem}


\section{The counted unit: Theorem 1 and its complete original proof}

\begin{theorem}\label{KANTO}
Let $X_1,\ldots,X_K$ be compact metric spaces, and let $\mu_k\in\mathcal{P}(X_k)$.
Let $c : X_1 \times \cdots \times X_K \to \mathbb{R}_+$ be continuous. Then the duality formula holds:
\begin{align*}
\mathsf{MOT}_c(\mu_1,\ldots,\mu_K)
=\inf_{\pi\in\Pi(\mu_1,\dots,\mu_K)} I(\pi)
\;=\;
\sup_{\boldsymbol{f}\in \mathcal{F}_c} J(\boldsymbol{f}).
\end{align*}
\end{theorem}

\subsubsection*{Proof of Theorem~\ref{KANTO}}
We recall from the Riesz representation theorem (see Theorem~\ref{RIEZ}) that the dual space \(E^{\ast}\) of $E=\mathcal{C}(X_1\times\ldots\times X_K)$ can be
identified isometrically with the space 
\(\mathcal{M}(X_{1}\times\cdots\times X_{K})\) of finite
signed Radon measures on the product, endowed with the total variation norm. Hence, to each \(L\in E^{\ast}\) corresponds a unique
\(\pi\in \mathcal{M}(X_{1}\times\cdots\times X_{K})\) such that
\begin{align*}
L(u)=\int_{X_{1}\times\cdots\times X_{K}} u\, d\pi.
\end{align*}
Our goal is to compute both sides of the duality relation provided by the Fenchel--Rockafellar theorem. The left-hand side is clearly
\begin{align*}
&\inf\left\{
\sum_{k=1}^{K} \int_{X_k} f_k\,d\mu_k
\;:\;
f_k\in C(X_k),\ 
\sum_{k=1}^{K} f_k(x_k)\ge -c(x_1,\ldots,x_K)
\right\}
\\
&\qquad =
-\sup\left\{
J(\boldsymbol g)
\;:\;
\boldsymbol g\in \mathcal F_c\cap \prod_{k=1}^{K} C(X_k)
\right\}.
\end{align*}

\paragraph{Legendre-Fenchel transform of \(\Theta\).}
For \(\pi\in \mathcal{M}(X_{1}\times\cdots\times X_{K})\),
\begin{align*}
\Theta^{\ast}(\pi)
&=
\sup_{u\in E}
\Big\{
\int u\, d\pi - \Theta(u)
\Big\}
=
\sup_{\substack{u\in E\\ u\ge -c}}
\int u\, d\pi.
\end{align*}
Hence
\begin{align*}
\Theta^{\ast}(-\pi)
=
\sup_{\substack{u\in E\\ u\le c}}
\int u\, d\pi.
\end{align*}
If \(\pi\) is not a nonnegative measure, then there exists
\(v\in \mathcal{C}_{b}(X_{1}\times\cdots\times X_{K})\), \(v\le 0\), such that
\(\int v\, d\pi>0\). Choosing \(u=t v\) and letting \(t\to +\infty\)
gives \(\Theta^{\ast}(-\pi)=+\infty\). If instead \(\pi\ge 0\), the supremum is attained at \(u=c\), and we obtain $\Theta^{\ast}(-\pi)=\int c\, d\pi$. Thus
\begin{align*}
\Theta^{\ast}(-\pi)
=
\begin{cases}
\displaystyle \int c\, d\pi, & \pi\in \mathcal{M}_{+}(X_{1}\times\cdots\times X_{K}),\\[0.3em]
+\infty, & \text{otherwise}.
\end{cases}
\end{align*}
\paragraph{Legendre-Fenchel transform of \(\Xi\).}
For any \(\pi\in \mathcal{M}(X_{1}\times\cdots\times X_{K})\),
\begin{align*}
\Xi^{\ast}(\pi)
=
\sup_{u\in \mathcal{C}_{b}(X_{1}\times\cdots\times X_{K})}
\Big(
\int u\, d\pi - \Xi(u)
\Big).
\end{align*}
Since \(\Xi(u)<\infty\) only when 
\(
u(x_{1},\ldots,x_{K})=\sum_{k=1}^{K} f_{k}(x_{k})
\),
this becomes
\begin{align*}
\Xi^{\ast}(\pi)
=
\sup_{(f_{1},\ldots,f_{K})}
\Bigg(
\int_{X_{1}\times\cdots\times X_{K}}
\sum_{k=1}^{K} f_{k}(x_{k})\, d\pi
-
\sum_{k=1}^{K} \int_{X_{k}} f_{k} \, d\mu_{k}
\Bigg).
\end{align*}
If
\begin{align*}
\int \sum_{k=1}^{K} f_{k}(x_{k})\, d\pi
=
\sum_{k=1}^{K} \int f_{k}\, d\mu_{k}
\quad\text{for all }(f_{1},\ldots,f_{K}),
\end{align*}
then \(\Xi^{\ast}(\pi)=0\). Otherwise, there exists a family \((f_{1},\ldots,f_{K})\) for which the
difference is nonzero. Scaling \(f_{k}\) by \(t\) (or \(-t\)) yields a
quantity that diverges to \(+\infty\), hence \(\Xi^{\ast}(\pi)=+\infty\). Thus
\begin{align*}
\Xi^{\ast}(\pi)
=
\begin{cases}
0, & \pi\in \Pi(\mu_{1},\ldots,\mu_{K}),\\[0.3em]
+\infty, & \text{otherwise.}
\end{cases}
\end{align*}

Both \(\Theta^{\ast}\) and \(\Xi^{\ast}\) are indicator functions:
\begin{align*}
\Theta^{\ast}(-\pi)=
\begin{cases}
\int c\, d\pi, & \pi\ge 0,\\[0.2em]
+\infty, & \text{otherwise},
\end{cases}
\qquad
\Xi^{\ast}(\pi)=
\begin{cases}
0, & \pi\in\Pi(\mu_{1},\ldots,\mu_{K}),\\[0.2em]
+\infty, & \text{otherwise}.
\end{cases}
\end{align*}
Therefore,
\begin{align*}
\max_{\pi\in E^{\ast}}
\big\{
-\Theta^{\ast}(-\pi)-\Xi^{\ast}(\pi)
\big\}
=
\max_{\pi\in\Pi(\mu_{1},\ldots,\mu_{K})}
\left(
-\int c\, d\pi
\right)
=
-\inf_{\pi\in\Pi(\mu_{1},\ldots,\mu_{K})}
\int c\, d\pi.
\end{align*}
Putting everything together and changing signs, we obtain
\begin{align*}
\inf_{\pi\in\Pi(\mu_{1},\ldots,\mu_{K})} I(\pi)
=
\sup_{\boldsymbol{f}\in \mathcal{F}_{c}\cap\prod_{k=1}^{K}\mathcal{C}_{b}(X_{k})}
J(\boldsymbol{f}).
\end{align*}
By Lemma~\ref{easyy}, this yields
\begin{align*}
\inf_{\pi\in\Pi(\mu_1,\ldots,\mu_K)} I(\pi)
=
\sup_{\boldsymbol f\in \mathcal F_c} J(\boldsymbol f),
\end{align*}
which completes the proof of Theorem~\ref{KANTO}.
\hfill\(\Box\)

\begin{thebibliography}{99}
\bibitem[19]{knapp2007} Anthony~W. Knapp. {\em Basic real analysis}. Cornerstones. Birkh\"{a}user Boston, Inc., Boston, MA, 2005. Along with a companion volume {\it Advanced real analysis}.
\end{thebibliography}
\end{document}
